\section{Recognition of presentations and compatible products}

All homomorphisms in this paper are group homomorphisms. A product of subgroups
means a product of sets until it has been shown to be a subgroup. No group is
assumed finite in this section.

\textbf{Theorem UA.Normal (normal kernel recognition).}
Let $\pi:P\to B$ and $\rho_i:P\to E_i$, for $i$ in a finite set $J$, be
homomorphisms. Suppose $K_i\unlhd P$ satisfy
\begin{align}
 \ker\pi&=\gen{K_i:i\in J},\label{eq:basekernel}\\
 \rho_i(K_j)&=1\quad(i\ne j),\label{eq:offdiag}\\
 \ker\rho_i\cap K_i&=1.\label{eq:faithful}
\end{align}
Then the homomorphism
\[
 P\longrightarrow B\times\prod_{i\in J}E_i,
 \qquad x\longmapsto (\pi(x),(\rho_i(x))_i)
\]
is injective. Moreover, the kernel in \eqref{eq:basekernel} is the internal
direct product of the $K_i$. For injectivity alone it suffices to replace
\eqref{eq:basekernel} by $\ker\pi\leq\gen{K_i:i\in J}$.

\begin{proof}
For $i\ne j$, an element of $K_i\cap K_j$ is killed by $\rho_i$ by
\eqref{eq:offdiag}, and is therefore trivial by \eqref{eq:faithful}.
Normality gives $[K_i,K_j]\leq K_i\cap K_j$, so distinct factors commute.
Their finite set product is consequently their generated subgroup. If
$\prod_i k_i=1$ with $k_i\in K_i$, applying $\rho_j$ gives
$\rho_j(k_j)=1$, whence $k_j=1$. Thus multiplication identifies the external
product with the generated subgroup. Finally, an element in the kernel of the
displayed homomorphism belongs to $\ker\pi$, hence has such a product
expression. Every coordinate vanishes, so every factor is trivial.
\end{proof}

\textbf{Corollary UA.RecognizeQuotient.}
Let $f:P\twoheadrightarrow H$ be a homomorphism, and let $H$ have a base
homomorphism $p:H\to B$ and coordinates $c_i:H\to E_i$. If
UA.Normal applies to $p f$, $c_i f$ and normal subgroups
$K_i\leq P$, then $f$ is an isomorphism.

\begin{proof}
The coordinate map factors through $f$, so injectivity of the coordinate map
implies injectivity of $f$. Surjectivity is the stated hypothesis.
\end{proof}

Normality is substantial. Let $P=D_8=\gen{u,v\mid u^2=v^2=(uv)^4=1}$,
let $B=1$, and take $K_1=\gen u$, $K_2=\gen v$. Its abelianization
$P\to C_2^2$ has coordinates faithfully detecting the two factors and killing
the other one. The factors generate $P$, but the map has the nontrivial kernel
$\gen{(uv)^2}$. Here neither factor is normal. Thus a correct local model and
coordinate detection do not themselves eliminate mixed relations.

The following specialization uses the full compatible product $H$ in
UA.SK and assumes that its support family is an antichain.

\textbf{Theorem UA.Recognition (complete presentation).}
In this setting, let $P=\gen{r_t:t\in I}$ be given by a
presentation. Suppose:
\begin{enumerate}
\item its relators hold on the marker tuples in the compatible product $H$;
\item the local relations of every $E_S$ hold on $(r_t)_{t\in S}$;
\item the image $K_S$ of $N_S$ under $E_S\to P$ is normal in $P$;
\item adjoining the relations $K_S=1$ for all $S$ gives exactly the marked
presentation of $B$.
\end{enumerate}
Then the marker map $P\to H$ is an isomorphism. In particular the $K_S$
form an internal direct product, all marked intersections are exact, and if
the input groups are finite then
\begin{equation}
 |P_L|=|B_L|\prod_{S\subseteq L}|N_S|.                         \label{eq:orders}
\end{equation}
The additional deletion conclusion of UA.SK.Delete transfers to $P$ under
its hypotheses. The same holds for the geometric conclusion of
UA.SK.Geometry under its hypotheses.

\begin{proof}
The first assumption and the generation assertion of UA.SK give a marked
surjection $f:P\twoheadrightarrow H$. The second gives $E_S\to P_S$;
its composite with $f$ and projection to $E_S$ is the identity on every
local marker, hence is the identity homomorphism. Thus each local map is
injective, without an order comparison. Under $f$, the image of $N_S$ is
its pure $S$-coordinate: the antichain makes every other intersection of
supports proper, where $q_T$ is injective. Consequently the coordinate maps
kill $K_S$ outside its own coordinate and are injective on $K_S$ in that
coordinate.

Normality makes $K=\gen{K_S:S\in\mathcal S}$ a normal subgroup. The last
assumption identifies the marked map $P/K\to B$ with an isomorphism, so
the base map $P\to B$ has kernel exactly $K$. UA.RecognizeQuotient
now proves that $f$ is an isomorphism. The remaining conclusions follow
from UA.SK and the explicitly indicated additional clauses.
\end{proof}

The last assumption is a presentation calculation. It does not assume the
isomorphism $P\cong H$ which is being proved. In practice it says that, after
killing the displayed local kernel words, the remaining relations are the
known base presentation. A pair of marked maps which are inverse on generators
also proves this condition. Normality likewise has a concrete local test: the
markers within $S$ already normalize $K_S$, and one must check the remaining
markers. It is enough that those remaining markers centralize a generating
set of $K_S$. Merely writing down the desired quotient does not supply this
check.


\textbf{Proposition UA.Mix (jointly faithful deletions).}
Let $G=\gen{r_i:i\in I}$ have every Boolean deletion $e_J$. Suppose
$I=L\cup M$, put $C=L\cap M$, and assume
$\ker e_L\cap\ker e_M=1$. Then
\[
 G\longrightarrow G_L\times_{G_C}G_M,\qquad
 g\longmapsto(e_L(g),e_M(g))
\]
is a marked isomorphism. The maps to $G_C$ on the right are the restrictions
of the corresponding deletions. No group is assumed finite.
\begin{proof}
UA.D gives $e_Le_M=e_C=e_Me_L$. Hence the displayed pair is compatible,
and the homomorphism's kernel is the stated trivial intersection.
For a compatible pair $(x,y)\in G_L\times_{G_C}G_M$, put
$c=e_C(x)=e_C(y)$. Since the two deletions fix their respective
parabolics, $e_M(x)=c=e_L(y)$ and both deletions fix $c$. Thus
\[
 e_L(xc^{-1}y)=xc^{-1}c=x,\qquad
 e_M(xc^{-1}y)=cc^{-1}y=y.
\]
Every compatible pair therefore has the explicit preimage $xc^{-1}y$.
The homomorphism is surjective as well as injective, with inverse given
by this formula. Its values on the markers give the stated marking.
\end{proof}

The kernel condition can be checked by coordinates. In a supported-kernel
realization with inherited deletions, suppose every base coordinate is
detected on at least one side, and every support is contained in $L$ or
$M$. An element killed by both deletions has trivial base coordinate. Each
pure supported-kernel coordinate is then detected on a side containing its
support, so all these coordinates are trivial too. This proves the kernel
condition in the path and two-branch applications below; it is not assumed
from a coincidence of finite orders.
